\chapter{第一集 \quad 秦汉}


\section{正文}


公元前223年，秦国和楚国的战事正陷入胶着，虽然楚王负刍已于一年前被秦军俘获，项燕和昌平军却仍在率军顽抗，即使骁勇如秦人，对上了楚人的坚韧，一时间也久攻难下。黑夫是秦军中的一名普通士兵，自参军以来，他和弟弟经常常并肩作战。然而这次两人却被分配到不同的战场，每一次分开，都可能是来不及告辞的生离死别。

安陆城原为楚国疆土，数十年前便已归秦国所有，如今这里居住着大量秦人，淮河边的战场距此不远，每一天都有家庭收到从军的亲人在战场上死去的消息，其余的家庭则在忧惧与思念中煎熬。因着两个弟弟都在从军，衷的一家也在等待中度日。这天终于有人带回了弟弟们的信，衷几乎是颤抖着展开，匆匆一眼便放下心来，赶紧向家人们转述信中内容。

黑夫和惊暂时无恙，前阵子他们分开作战，现在又见面了。黑夫和惊此时正在淮阳攻打反城已久，前路未卜，伤未可知。他们希望母亲寄来的钱一定不要少，还反复叮嘱衷，务必留意官府的军功授爵文书。

信末一行行罗列着对熟识亲友们的絮絮问候，最后一行是惊对妻女的关怀，他并嘱托妻子勉励照顾家中老人。身处生死瞬息的战场，这些琐碎的牵挂才是战士们心里最坚强的铠甲。安陆城外的连绵丘垄，多是城中居民们的死后安息之所。黑夫和惊曾经牵挂的亲人就葬在西郊一片，后来被称为“睡虎地”的坡地上。

20多个世纪后，考古工作者在其中的四号墓里发现了当年黑夫和惊寄出的信，信写在木牍上，木牍被珍重放置在木椁头箱中，收件人衷或许就是这座墓的墓主。

彼时秦国像黑夫和惊这样的平民，只有履险阻，冒矢石，以性命相系，方能换得军功之爵，以此改善个人与家庭的境况，或许他们最终还是生死暌违，家书便成为了生随死葬的珍藏。从辽东至陇西，自长城至南海，数百年间干戈不止，纷争不休。秦始皇陵兵马俑那每一张都迥异的面目，是曾披肩之锐的赳赳雄狮，也是曾被生离死别磋磨的血肉躯体。他们以及六国故地那更多未曾留名的战士，艰险岁月里为己即是为家，为国亦是为家。与家人间的牵挂，就是暗夜里明亮的光。

黑夫与惊写下家书的两年后，即公元前221年，世人期盼已久的和平局面崭然显露，横扫六合的秦军终能暂时放下手中的武器，看看那曾为之浴血的秀丽山川。

一座座曾经倾颓的城邑内外重新升起了连绵不断的炊烟，数百万平方公里的山河被纳入以中央政府为枢纽的郡县网络中，前所未有的辽阔王朝沐浴在初升的朝晖之下，巍巍伫立。

在连年的征战中，秦朝君臣早已认识到，东周列国之间在制度上无法兼容的差异，乃是攻扦不息的根源之一。度量衡、车轨、钱币、法律和文字，因为涉及到行政、手工业、商业和社会治理等重要领域，被迅速选定成为秦朝建立后需要统一的关键内容。

曾经的六国故地和秦都咸阳所在的关中，都曾经发现大量镶嵌秦诏版的铜权和铁权。权是称量重量的秤砣。权上的诏版写道，皇帝并天下后，百姓安定，令丞相隗状、王绾规范度量衡。凡不一致者，皆需统一。这道统一度量衡的诏书也被模印或铭刻在官定的度器和量器上，流布天下。

公元前219年，秦始皇在仲春之月开始了他的冬巡。此后八年间，他一度封禅，五次冬巡。他登峰山、攀泰山，上琅琊，经之罘，过东观，临碣石，至会稽，巡行之处皆刊石勒铭。

这七处石刻，由于风雨剥蚀，人事相侵，仅有琅琊、泰山二石尚存残玦，峰山刻石尚于摹刻本。刻石文辞以李斯新改定的小篆向关东宣示：乃今皇帝壹家天下，兵不复起，灾害灭除，黔守康定，利泽长久。秦要统治的是文字语言存在极大差异的四海八方，秦朝的文字改革使得汉字形体简化，部首和笔数固定，文字使用的标准得到统一。

湖南里耶古城遗址出土的这枚秦代木方，由六个残片拼接而成，以较为拙朴的篆书记录了其时文字的变更。“酉”如故更“酒”，将酒从酉分出，凡酒之意均用“酒”字，不再写作“酉”；“卿”如故更“鄉”，意思是公卿之“卿”仍然照旧，而记入鄉里之“鄉”，统一更用“鄉”字。

经过规范的文字书写和传播起来更有效率，使中央的政令能通过公文畅通地下达至各个郡县，也使更多人能够通过文字交流获得知识和文化的传承。书同文成为了中华文明赓续几千年而始终文化一体的奥秘。

渤海湾畔至今仍存多处秦代行宫遗址。在辽宁绥中万家镇南边，就发现了六组相互关联的秦汉建筑遗址，其中三组皆在海岸，朝向海中的“姜女石”呈合抱之势。这或许就是当年秦始皇、汉武帝曾登临过的“碣石”。石碑地建筑群可能就是秦汉时的“碣石宫”。

曲尺形宫墙之内，建筑高低错落，疏密相间，依稀可见昔时廊道鳗回，院落镶嵌的格局。洪波高涌，碣石为阙，此为秦之国门。临碣石而东望，日升月落，天旋海沸，终至陆地之极，身后是万里河山。

这是秦始皇东巡时，臣属们镌刻在琅琊台石碑上的文辞，残石饱经沧桑，铭文也大多漫浊不清，但内容依然震撼人心。

六合之内，皇帝之土，西涉流沙，南尽北户，东有东海，北过大夏，人迹所至，无不臣者。但天下并不太平，北方边境地平线上争端频仍，匈奴的铁蹄声一旦响起，劫掠与杀戮就会到来。

30万秦军正等待主帅号令，蒙恬已记不清自己经历了多少场战役，只知道在争夺天下时，他是秦王的长策，四海平定时，他就是皇帝的敲扑。面对匈奴滋扰，蒙恬率军却敌，把秦朝的边境线向北推进700余里，直抵阴山脚下。为了在外敌再度来犯时能够迅速地预警并组织反击，蒙恬受命将原先秦、赵、燕三国的长城与由洮水而来的长城相连，沿途起城邑修亭障，随后领兵驻守上郡十余年。

秦始皇35年（公元前212年），蒙恬率军修建了直道。这条长约743公里、宽达50米的道路，以几近直线的走向连接了关中与北方草原地区。经行于山地和丘陵时，削平山脊，填平山谷；经行于沙地和草原时，以夯打或掺石的方式加固路面；经行山侧河畔时，设置护坡和排水沟。这不仅是皇帝巡行北方的“高速公路”，更是支援长城军事防务的运输线。

长城如盾，直道就是持盾的手臂，它们共同铸就了延袤万余里的坚固藩篱。百余年后，史官司马迁行经直道时，忍不住为这项暂山埋谷的工程惊叹。然而，他也指出长城和直道费人力财力无数，虽与王朝的防守有利，却施之过急，未留给刚刚从战乱中安定下来的众生以喘息之机。尽管赢政在“始皇帝”的名号里寄寓着他千秋万世的构想，却仅二世而亡。

直到汉才带来了较长时间的大部分疆域的和平。作为汉朝的开启者，刘邦最初选择的都城是洛阳，他倚中的大臣都是关东人，无不希望永久地定都于此，只有齐人娄静和留侯张良力主迁回关中。虽然战事以息，局势却仍不稳定。关中坐拥天险，不易受敌人威胁，且土地膏腴，交通便利，未来可容纳更多的人口。此外，还可以割断关东功臣集团们与乡土间的利益关系。公元前202年，刘邦带着他的朝廷西迁后暂居于栎阳，这里是秦人旧都，商鞅在此变法，秦的壮大也正是从这里加速。

两年后，在50多公里以外的渭河南岸，由秦朝渭南离宫之一兴乐宫改建而成的长乐宫竣工，意味着新王朝都城的诞生，长治久安的愿望被寄托在这座新都长安的名称里。

未央宫紧闭的宫门一年后轰然洞开。刘邦为这座新宫室的华丽所震惊，即便他见过秦朝巍峨的咸阳宫，其宏伟也远不及眼前的未央宫。他想起了秦王的教训，想起了刚刚剿灭的叛乱，作势怒斥主持工程的萧何，萧何俯首答道“夫天子四海为家，非壮丽无以重威，且无令后世有以加也。”刘邦已领略了威家四海的滋味，这话令他大悦，但清醒之后仍心有惕惕。此后至死也将主要的政治活动放在长乐宫中。

西汉都城长安是当时世界上规模最大的城市之一。龙首原上，长乐、未央两宫东西并峙，于高处制掣全城。桂宫、北宫、明光宫横列两宫之北。长安有东西二市，八街九陌，12城门160闾里。

渭水之北，帝陵有九，五陵邑拱卫京师，东郊多丘墓，东南郊有文、宣二帝之陵，南郊为礼制建筑，西郊则是离宫苑囿。

200年间多少政治风云由此升起，多少人事兴亡在此兴起，个个后席卷着汉朝疆域的每个角落。使臣陆贾从长安出发，携带着汉朝预备封授给南越王的印绶前往番禺。

南越王赵佗本是秦朝派往岭南的一名地方官员，秦末趁乱据五岭以南自立为王。刘邦平定岭北后，不想再冒着瘴疫的危险翻山越岭，尤其越地丛林深密，水道纵横，并不是汉军熟悉的作战环境。他希望用一种和平的方式，令赵佗主动归顺。曾随高祖征战的陆贾，对一路的旧关防再熟悉不过。如今关防已撤，商贾周流，天下繁华气象初现。

过了五岭便是南越，陆贾一行弃车马，换舟楫顺江而下，鹧鸪声里，两岸的木棉树穿古榕而出，绿水明波荡漾。

番禺城中央的宫苑，自然不及汉家都城宏丽，却处处可见效仿之意。赵佗免冠露髻在大殿之上交脚而坐，倨傲地望向汉家使臣。这在礼仪严明的中原人看来，他分明以走成人自居，摆出了割裂的姿态。陆贾款款陈词，既点出赵佗本为中原人的身份，又在对时事的条分缕剖中暗藏了威慑与怀柔。

听毕，南越王肃然起身，重新庄重地按中原礼仪跪坐相待。赵佗先征询得到陆贾认为他比汉家丞相们更贤能的答案后，又在追问中将比诏的对象换成汉家皇帝。陆贾看出南越王心思，并未直接评价二人才干，却转而分析汉越之间的实力差异。身在岭南，无法与中原王朝抗衡，赵佗如何不明白其中因由，如今借陆贾之口说明白，倒也让多年心事了结。虽然股肱大臣们多是中原人，但赵佗的心里是寂寞的，无人能了解他欲振翅于寰宇的鸿鹄之心，越中无人可共语。陆贾这一介使臣的才干胆识，却令他引以为知交。赵佗已经可以想见汉家朝廷的英才济济。

此后赵佗常常邀陆贾饮酒谈天，有时候宴会设在华音宫畔的曲流石渠旁，有时在波光粼粼的蕃池旁。宫苑中，梅花鹿撷食着灌木的嫩芽，龟鳖从渠畔的斜坡爬上岸享受阳光，宫人兢兢业业地清点每颗从北边移植来的枣树上结了多少颗果实。如此数月，赵佗在听取陆贾对天下大势的剖析后，表示从此臣服汉室，并馈赠陆贾宝物，送其北归。

昔日的南越宫苑虽毁，周遭的繁华却2000年未息，其遗址便沉睡在广州的闹市之下。曾经流水蜿蜒的石渠；曾经天光云影的蕃池，至今还回放着历史每一次的柳暗花明，峰回路转。

一件陶提简盖残片上的戳印暗示了赵佗的心境——华音宫。汉初的君臣常常讨论，为何秦统一了疆域开创了前所未有的辽阔版图，却短祚覆亡？在反思中他们逐渐达成共识，柔化统一之策略，厚积统一之基础，化育民意于无形，静水自可流深。

在数代君臣的治理下因长久战乱而凋敝的民生逐渐恢复，同时恢复的还有那曾因为秦而失去的对于“大一统”的信心。武帝接过的是一个初具气象的王朝，此时董仲舒挟数世儒生积累的智慧而来，提出依托儒家思想构建大一统下的国家认同。这显然触动了汉武帝的所思与所愿。他看到了让六合同风、四方共俗的可能性；看到了不再依靠血缘关系治理社会的可能性。南昌市新建区的郭墩山下，沉睡着刘贺夫妇及其子孙，他的祖母是生前宠冠后宫，死后与武帝合葬的李夫人。刘贺继承了父亲的昌邑王位，昭帝崩逝后，迅速被霍光等定为新帝人选。然而在帝位上仅仅待了27天，刘贺便成为关中群臣攻讦的目标，背负无数罪名被褫夺王位，只得以列侯之身终老海昏。

海昏侯墓内建有400平方米的木构椁室，主墓室的设计模仿刘贺生前的真实生活场景，事死如事生。墓中随葬的大量黄金和精美器物可以窥见当时国家的强盛富足。在主椁室的西室内，摆放着莲枝灯、博山炉和漆案。这面漆木衣镜屏陈列在床榻旁最显眼的位置，镜屏是刘贺生前所用陪伴他走过人生大半，镜背表面绘着孔子和他的五个弟子，这是目前所见年代最早的孔子图像。人像两侧以墨书抄录着人物生平及言行：被任用时就施展抱负，不被任用就藏身自好，只有我和你才能这样吧。

当年汉宣帝将刘贺封至南昌做海昏侯，实质上是要他离开富裕的山东到偏远的“南藩”豫章就国还令他不得回长安祭祀祖先。

关山难越，谁悲失路之人，萍水相逢尽是他乡之客。情难自己时，海昏侯乘流东望，辄愤慨而还。

刘贺墓内随葬大量文书，除签牌、奏牍等外，以规整抄录的儒家经典居多，唯有一版摘抄了《论语》断篇的木犊，字体率性随意与别不同。在废黜之后被时时监视的日子里，刘贺无以纾解，唯有寄情于经典之间。他亲笔抄写下了这样的字句“子曰：吾者知乎哉？毋知，有鄙夫问乎吾，空空如也。叩其两端而竭。”人生多疑惑而少智慧，只得穷极叩问，竭尽求索，或许刘贺就在圣贤的教诲间寻找着安身立命的心境。

武帝推崇儒家像刘贺这样的王侯往往能以大儒、博士为师，但对于普天下的学子来说，负笈求学于都城，才是他们的夙愿和殊荣。自武帝立太学后办学规模日渐攀升，到质帝时，从各地前往洛阳求学的学子已达到3万人之多。太学中黄舍上千，鳞次栉比。灵帝召集当时鸿儒正订六经，刻石立于太学。在没有印刷术的年代这“熹平石经”就是学子们的教科书，前来观视与临摹石经的车骑络绎不绝，甚至一度堵塞了太学门口的街道，秦完成了形式上的统一。

汉则使统一的理念嵌入到思想和文化的基因里，以秦朝的文字统一为根系，两汉以太学为参天巨木之干，将治国的理念与秩序寄寓在教育之中，快速抽生出荣茂的枝叶，伸向万里河山的各个角落。即使在当时的偏远之地，都曾有同样熟悉的诵读声，普天下的莘莘学子，从高门绣户的列侯到边陲关隘的无名小卒，西域精绝城里的当地人也都曾捧着同样的《仓颉篇》识文习字，念着同样的冬寒夏暑，玄气阴阳。汉兼天下，海内并廁。他们都将成长为治理这片辽阔疆土的力量。

在注重文治的同时，国家安全仍然是汉朝最为关注的议题。武帝在位时国内人口充盈、物资丰饶，诸侯国也再不能掣肘，海内为一，兵马整肃，但匈奴在这期间也发展壮大，对北境的安全造成了威胁。

公元前138年，汉武帝公开招募西行联络月氏、大渊和乌孙以抗击匈奴的使者。27岁的汉中青年张骞毅然走上“凿空”西域之途，这一去就是整整13年。

从长安出发时，车马粼粼百余人，回程路上只有两人，13年间有利益相诱，有武力相胁，温柔乡、囹圄境，却从未丢失身为汉使的旌节，从未磨灭此行的初心。

公元前130年，年轻的车骑将军卫青直捣龙城，获得了汉朝对战匈奴的初次胜利。公元前123年，嫖姚校尉霍去病率轻骑直入大漠。两年后，又在祁连山奇袭匈奴的浑邪王、休屠王。

公元前119年卫青、霍去病远征漠北，令单于遁逃，封狼居胥。八年后，汉王朝先后设酒泉、张掖、敦煌、武威四郡，牢牢掌控河西走廊。此后匈奴虽然还有零星骚扰却已元气大伤，不再成为气候。尽管如此，庞大的戍边队伍仍在长城沿线屯驻，兵马不懈卫护着汉朝北境的绝对安全。

转也是其中的一员。秋风渐凉，这日正值休假，同袍们有的凑在一起玩六博戏，有的去与随军的家人团聚，有的约了朋友聚会饮酒……转听闻此前受托去家中探视的友人幼卿，此时已来到相距不远的肩水都仓，很想赶去见上一面，却因公事无法脱身，只得修书一封托人转交。他请幼卿往家中捎话“老人平安便是离乡远行者之幸。天气凉了，千万要记得及时添衣，善进饮食。”正要搁笔时想起秋天是匈奴人经常进犯的季节，赶紧又添上几句嘱咐“幼卿你独自行走于关外，此身只一并无二，千万当心。”

转和幼卿如此平凡，他们漫长的人生我们只能窥见这只字片语的时光，再不知后事如何。转写给幼卿的那封信也一裂为二，散落坞墙之下，长埋黄沙之中。

1973年的夏天，大西北的戈壁上仍旧延续着年复一年如期而至的高温和风沙。昔年的居延曾有“弱水流沙”之称，此时“弱水”已更名为额济纳河，流沙依旧无休止。

从7月到9月，甘肃居延考古队都在肩水发掘，与秦时无两的明月照在汉时的关址上，连流沙和相思也并无两样。发掘常常被风沙中断，停了之后又再重来，如此历经将近三个月，汉代河西的这一座金关终于再度重现。

37个探方中11577枚汉简记录着2000多年前的文书与人事往来。考古工作者耐心地将简牍一枚枚清洗、分类、缀合、释读。在这个过程中，30号探方和26号探方中的两片断简被合而为一。碎了的信重新还原，我们才成为时隔二十个世纪的读信人。

为汉朝卫护北部边境国家安全的不仅有使者、将士也有和亲的公主们。乌孙是西域诸部族中较为强大的一个，随畜逐水草而居。张骞向武帝献策，以财物相诱送公主和亲，说服乌孙与大宛、月氏形成夹击匈奴之势。

解忧公主是楚王刘戊的孙女，她的爷爷在景帝时参与了一场史称“七国之乱”的反叛行动，最终兵败自杀。或许因为皇室宗亲和罪臣后嗣的双重身份，她成为汉朝派往乌孙的第二位公主。此前不久，江都王女刘细君刚刚死在异乡，细君公主生前所作的悲歌早已传遍长安“居常土思兮心内伤，愿为黄鹄兮归故乡。”

与草长莺飞的故乡迥异的风物与习俗，是武帝诏书中“欲与乌孙共灭胡”的重托也无法冲淡的忧愁。在政治漩涡中成长的解忧有着与细君公主不同的坚韧心性。她带着王朝的托付，一意西行陪她同往的还有聪慧的侍者冯燎。

从长安向西他们第一次看见天地间如此空荡的景象，四面八方的风无所阻滞地在莽莽平野间来往。第一次看见尘世间如此壮阔的奇观，雪山、沙漠、戈壁、龙堆以及镶嵌其间的湖泊与绿洲。

到乌孙后，解忧嫁与国王军须靡，冯燎则为右大将之妻。他们居毡帐，饮酪浆，白天听着群马嘶鸣，夜晚枕着猎猎风声。对故乡的思念此刻被化作对故土的使命，开始了他们在西域长达数十年的斡旋。

君须靡死后，解忧遵照乌孙的习俗嫁给了他的堂弟翁归靡。公主在内润物无声的影响着乌孙的政局。她的新丈夫不再游移于汉匈之间，而选择了一心归附于汉。冯夫人在外，持汉节为公主行赏赐，邀迎西域诸部族之心。

解忧在乌孙40余年的经营眼看即将换来长久的安宁，翁归靡却猝然薨逝。

60岁的解忧不得不嫁给了她的第三任丈夫泥靡。这是她第一任丈夫与匈奴人的儿子，对汉朝并无好感，她处境艰难，一度生死悬于一线。事情的转机发生在泥靡被杀之后，在冯燎的帮助下，在汉庭的支持下，他的儿子元贵靡最终成为了乌孙的大昆弥。一度屯积于敦煌待命的汉军得以不战而还。

这一年是甘露二年（公元前52年）解忧已经年近七旬，她尚记得自己来时明眸皓齿，风华正茂，如今遥望家乡时，她的眼神却越来越模糊，她怕再不回去便忘记故土的模样。她一生里大部分的时光都在这伊犁河畔度过，为了故国所赋予的使命，变了语言，改了习惯，默默地适应着这里的一切。

她的三任丈夫，她的儿女都在这里，但是她不能忘记自己的来处，那灼灼的桃花，那苒苒的桑梓，那方正的文字，那绮丽的丝绸，这也是千百年后仍然嵌在中国文化里的基因——对故乡的眷念始终深深镌刻在骨血里。

从玉门到敦煌一路皆是茫茫戈壁，汉代的马甲骑联翩，也曾由此经行。在距敦煌市区还有60公里的柳格高速南边，悬泉置伫立在连绵而荒凉的山脉之下，这是敦煌郡效谷县属下的一处邮驿机构。日常在此工作的吏卒有30余名，负责传递官方文书、军情急报，接待往来的各级官员和各方使者。

啬夫弘是悬泉置存在的前后数百年时间里任职最长的一个。从宣帝元康三年到元帝初元四年的十八年间，他几乎未曾离开过这里。他为都护、使者、公主、将军、列侯们传过无数封信，准备过无数次行程所需的马匹与粮草。他仍清楚记得12年前他刚到悬泉置不久，长罗侯常惠护送少主相夫来到敦煌，不久后又因和亲不成而折返，之后便听说了解忧公主和冯夫人那惊心动魄的经历。

公元前58年，他看见丞相史李尊护送复原的戍卒返回原籍，在归乡的队伍中总有招车上载着的一具具棺木，那是不幸死在异乡的戍卒们。曾经，新征的年轻人从悬泉置经过，去往敦煌等地服役。如今他们也经这里返回各自的家乡。河东、南阳、颍川、东郡、魏郡、淮阳国，目送那些曾卫护过一方安宁的人们离开。啬夫弘仔细录下公文中戍卒们的原籍。

甘露二年（公元前52年）二月十二日傍晚，平望驿骑带来了长罗侯和解忧公主的书信，他们的书信一向与西域局势有关。啬夫弘未敢怠慢，将之交托给驿骑朱定，让他连夜赶往万年驿。

一年后啬夫弘收到了下行的公文，让他为从乌孙归汉的解忧公主一行准备车驾，他这才知道曾送出的那封信，正是公主归家的恳求。那年冬天他看见年老的解忧公主和她的三个儿孙匆忙上前迎接。

从赤古城到悬泉置是1273公里，从悬泉置到长安城是1372公里，这位离家50年的大汉公主已在遥远的西域看过了将近2万轮的日升月落，古稀之年的她终于踏上了归途。岁月轮转，烽火已熄，古道上音尘渺绝，戍卒的书信埋在黄沙下，王侯的宫苑隐在高楼间，而那些城邑、道路、边关、堰渠曾见证着秦汉时期每一个鲜活的人，如何在这片山河湖海间努力经营着自己的一生，慢慢垒就矗立历史长河中的巍巍王朝。

他们曾眼望星辰；他们曾背负长天；他们曾跋山涉水；他们曾生死离别；他们曾在不知觉间成为两千年后今日中国的奠基。2000多年前长安城中的汉武帝望向九州华夏的舆图，他决意为汉王朝打开连接外部世界的通衢，沿着士兵、公主、使臣们曾经往来的河西走廊，越过葱岭向西而行；沿着商贾、船队曾经航行的线路，穿过大海向西而行。海陆并行的通途寄寓着强大的国家对外文化经贸交流的渴望，也孕育着中国大地此后对外交通的格局。

位于山东临沂的北寨汉画像石墓，刻画着一幅充满日常生活气息的乐舞百戏图。跳丸、飞剑、走索、寻橦……文献中记载的民间杂技鲜活的展示在人们面前，更特别的当属经丝绸之路传来的“鱼龙曼衍”之戏。据说这种戏使用西域幻术能够呈现出鱼龙变化的神异场景，不知不觉间外来文化沿着河西走廊流入生机勃勃、兼容并蓄的东方。走入寻常百姓家，而曾经横亘在东西方之间的山海重重已成通途。

统一而不断强大的中国，从秦汉时代启程，通过强有力的中央政府建构起多元一体的文明持续发展至今。

独特的中华文明如何形成？秦汉之前是否有更早的中国雏形？三皇五帝的古史系统是否可靠？中国文明5000年是否可信？让我们追随考古学家的脚步，作万年的时光旅程。

看我们的文明在摇篮里的孩提模样，被她稚嫩的笑容打动，看她在辽阔的山川大地上蓬勃成长，在5000多年前形成一个可以称作中国的共同体，长成为真正的文明；看她经历龙山时代500年激荡整合，将共同体凝聚为政体，为夏商王朝的诞生欢欣鼓舞；看她究天人之际，由尊帝到尚德，由编户至齐民，完成政体、社会乃至华夏文明的整合。

我们自秦汉溯源，探寻何以中国，贯通五千年文明血脉。我们经秦汉传续，虽历五千年而不老，中国依旧如旭日朝阳。


\section{典故}

\begin{table}[htbp]
\centering
\small
\renewcommand{\arraystretch}{1.2}
\caption{\textbf{第一集“秦汉”典故与文化知识}}
\begin{tabularx}{\textwidth}{@{}c p{3cm} L@{}}
\toprule
%每列的表头内容居中
\multicolumn{1}{c}{\textbf{序号}} & \multicolumn{1}{c}{\textbf{典故/文化名词}} & \multicolumn{1}{c}{\textbf{文中相关语句（或出处）}} \\
\midrule
1 & 黑夫与惊 & “黑夫是秦军中的一名普通士兵……黑夫和惊暂时无恙” \\
2 & 睡虎地秦墓 & “后来被称为‘睡虎地’的坡地上……四号墓里发现了当年黑夫和惊寄出的信” \\
3 & 秦始皇封禅（shàn） & “他一度封禅，五次冬巡” \\
4 & 李斯小篆（zhuàn） & “刻石文辞以李斯新改定的小篆” \\
5 & 书同文 & “书同文成为了中华文明赓（gēng）续几千年而始终文化一体的奥秘” \\
6 & 碣（jié）石宫 & “石碑地建筑群可能就是秦汉时的‘碣石宫’” \\
7 & 蒙恬（tián）修长城与直道 & “蒙恬受命将原先秦、赵、燕三国的长城相连”“蒙恬率军修建了直道” \\
8 & 司马迁叹直道 & “百余年后，史官司马迁行经直道时，忍不住为这项堑（qiàn）山堙（yīn）谷的工程惊叹” \\
9 & 萧何建未央宫 & “夫天子四海为家，非壮丽无以重威” \\
10 & 陆贾出使南越 & “陆贾款款陈词……赵佗肃然起身，重新庄重地按中原礼仪跪坐相待” \\
11 & 赵佗（tuó）与南越国 & “南越王赵佗本是秦朝派往岭南的一名地方官员，秦末趁乱据五岭以南自立为王” \\
12 & 华音宫 & “一件陶提筒盖残片上的戳（chuō）印暗示了赵佗的心境——华音宫” \\
13 & 董仲舒 & “董仲舒挟（xié）数世儒生积累的智慧而来，提出依托儒家思想构建大一统下的国家认同” \\
14 & 海昏侯刘贺 & “南昌市新建区的墎（guō）墩山下，沉睡着刘贺夫妇及其子孙” \\
15 & 关山难越，萍水相逢 & “关山难越，谁悲失路之人？萍水相逢，尽是他乡之客” \\
16 & 熹（xī）平石经 & “灵帝召集当时鸿儒正订六经，刻石立于太学” \\
17 & 太学 & “自武帝立太学后办学规模日渐攀升……到质帝时已达3万人” \\
18 & 张骞（qiān）凿（záo）空 & “27岁的汉中青年张骞毅然走上‘凿空’西域之途” \\
19 & 卫青、霍去病 & “卫青直捣（dǎo）龙城……霍去病封狼居胥（xū）” \\
20 & 封狼居胥（xū） & “令单于遁逃，封狼居胥” \\
21 & 解忧公主 & “解忧公主是楚王刘戊（wù）的孙女……一意西行” \\
22 & 冯嫽（liáo） & “聪慧的侍者冯嫽” \\
23 & 悬泉置 & “悬泉置伫（zhù）立在连绵而荒凉的山脉之下……啬（sè）夫弘是任职最长的一个” \\
24 & 居延汉简 & “37个探方中11577枚汉简记录着2000多年前的文书与人事往来” \\
25 & 鱼龙曼衍（yǎn） & “经丝绸之路传来的‘鱼龙曼衍’之戏” \\
26 & 龙山时代 & “经历龙山时代500年激荡整合” \\
27 & 究天人之际 & “看她究天人之际，由尊帝到尚德，由编户至齐民” \\
\bottomrule
\end{tabularx}
\label{tab:heyizhongguo_fixed}
\end{table}


\section{经典词语}

\begin{table}[htbp]
\centering
\small
\renewcommand{\arraystretch}{1.2}
\caption{\textbf{第一集“秦汉”精彩成语与词语（注音+解释+例句）}}
\begin{tabularx}{\textwidth}{@{}c p{2.6cm} p{2.6cm} X@{}}
\toprule
\multicolumn{1}{c}{\textbf{序号}} & \multicolumn{1}{c}{\textbf{成语/词语}} & \multicolumn{1}{c}{\textbf{名词解释}} & \multicolumn{1}{c}{\textbf{示例句}} \\
\midrule
1 & 胶着（jiāo zhuó） & 比喻相持不下或工作不能进展，如胶黏着。 & “敌人西侵的部队在长江的南北两岸都呈着胶着状态了。”——郭沫若《羽书集·致华南友人们》 \\
2 & 生离死别（shēng lí sǐ bié） & 很难再见的离别或永久的离别。 & “独有王夫人和宝钗娘儿两个倒像生离死别的一般，那眼泪也不知从那里来的，直流下来，几乎失声哭出。”——曹雪芹《红楼梦》第一百十九回 \\
3 & 忧心忡忡（yōu xīn chōng chōng） & 形容忧虑不安的样子。 & “忧心忡忡，待旦而入。”——韩愈《送李愿归盘谷序》 \\
4 & 絮絮叨叨（xù xù dāo dāo） & 形容说话啰嗦，连续不断。 & “那婆子絮絮叨叨，只管说些闲话。”——吴敬梓《儒林外史》第二十回 \\
5 & 披坚执锐（pī jiān zhí ruì） & 穿着坚固的铠甲，拿着锋利的武器。指奔赴战场作战。 & “披坚执锐，临难不顾，身先士卒。”——陈寿《三国志·魏书·曹彰传》 \\
6 & 生死暌违（shēng sǐ kuí wéi） & 生与死分离，不得相见。暌违：分离，隔离。 & “新知虽已乐，旧爱尽暌违。”——何逊《赠诸游旧》 \\
7 & 干戈不息（gān gē bù xī） & 战争不停息。干戈：古代兵器，借指战争。 & “干戈不息，征徭未休。”——白居易《与元九书》 \\
8 & 赳赳武夫（jiū jiū wǔ fū） & 威武雄壮的武士。 & “赳赳武夫，公侯干城。”——《诗经·周南·兔罝》 \\
9 & 李代桃僵（lǐ dài táo jiāng） & 以此代彼，代人受过。原比喻兄弟相爱相助。 & “桃生露井上，李树生桃旁。虫来啮桃根，李树代桃僵。”——《乐府诗集·鸡鸣篇》 \\
10 & 鳞次栉比（lín cì zhì bǐ） & 像鱼鳞和梳子齿那样有次序地排列着。多用来形容房屋或船只等排列得很密很整齐。 & “市井街坊鳞次比，鲛绡泉贝积如山。”——北宋·李邴《咏宋代泉州海外交通》 \\
11 & 势如破竹（shì rú pò zhú） & 比喻节节胜利，毫无阻碍。 & “今兵威已振，譬如破竹，数节之后，皆迎刃而解。”——房玄龄等《晋书·杜预传》 \\
12 & 背水一战（bèi shuǐ yī zhàn） & 比喻决一死战。 & “信乃使万人先行，出，背水陈。赵军望见而大笑。”——司马迁《史记·淮阴侯列传》 \\
13 & 直捣黄龙（zhí dǎo huáng lóng） & 比喻直捣敌人巢穴。 & “直捣黄龙府，与诸君痛饮耳。”——《宋史·岳飞传》 \\
14 & 封狼居胥（fēng láng jū xū） & 比喻建立显赫功勋。 & “骠骑将军去病率师躬将所获荤允之士，登临瀚海，封狼居胥山。”——班固《汉书·霍去病传》 \\
15 & 筚路蓝缕（bì lù lán lǚ） & 形容创业艰苦。 & “筚路蓝缕，以启山林。”——左丘明《左传·宣公十二年》 \\
16 & 家喻户晓（jiā yù hù xiǎo） & 家家户户都知道。形容人所共知。 & “使其家喻户晓，则人人自危。”——欧阳修《乞罢刽子配隶人户札子》 \\
17 & 前所未有（qián suǒ wèi yǒu） & 历史上从来没有过。 & “而中国政体，实为前所未有之奇变。”——梁启超《新民说》 \\
18 & 浴血奋战（yù xuè fèn zhàn） & 形容顽强地拼死战斗。 & 现代常用成语，多见于抗战文学。例：“八路军在阵地上浴血奋战，终于守住了防线。” \\
19 & 惊心动魄（jīng xīn dòng pò） & 形容使人感受很深，震动很大。 & “其文则惊心动魄，一字千金。”——钟嵘《诗品》卷上 \\
20 & 欣欣向荣（xīn xīn xiàng róng） & 比喻事业蓬勃发展，兴旺昌盛。 & “木欣欣以向荣，泉涓涓而始流。”——陶渊明《归去来兮辞》 \\
\bottomrule
\end{tabularx}
\label{tab:excellent_words}
\end{table}